\chapter{Objetivos}
\section{Objetivo Principal}
\begin{itemize}
	\item Analizar, diseñar e implementar un robot móvil para asistir en la 
		recuperación de pacientes con inmovilidad en los brazos.
\end{itemize}
\section{Objetivos Específicos}
El objetivo principal es claro y muestra perfectamente las intensiones del presente
trabajo, existen varios objetivos específicos, sin embargo he decidido que los 
objetivos específicos sean mutables a lo largo de la investigación y desarrollo,
la razón no involucra ir acomodando las expectativas a mi beneficio si no 
exactamente lo contrario, al momento de iniciar no tengo todos los detalles y es
ingenuo pensar que con el poco conocimiento que tengo sobre el problema puedo 
realizar la mejor solución, el lector quizá pensará que una investigación 
previa de material externo dará la información suficiente, pienso que esto no 
es así, la mejor forma de alcanzar todos los requisitos es obtener información del
problema mismo, de esta forma es él quien nos guía sobre el desarrollo,
derivando en un procesos de
constante iteración y realimentación.

Los objetivos a continuación son los que considero importantes en este instante 
(en el inicio del presente trabajo), pero sé y adelanto que irán cambiando, pues
el propio problema nos irá indicando que es lo que requiere para su solución.
\begin{itemize}
	\item Para el paciente.
		\begin{enumerate}
			\item Mover el robot de un punto A a B.
			\item Potencia suficiente para soportar la carga.
			\item Movimiento save.
			\item Rotación libre.
		\end{enumerate}
	\item Para el médico.
		\begin{enumerate}
			\item Interfaz de trayectorias.
		\end{enumerate}
	\item Para el desarrollador.
		\begin{enumerate}
			\item showoff.
		\end{enumerate}
\end{itemize}
